%% kapa: exam=1 round=19 year=2008 period=1 form=A subject=civil_law subject_ko=민법 seq=22 number=22
%% source: Q-Net 감정평가사 기출문제 (cst00309) artlSeq=1001962 file=extracted/1008/1001962/2008년 제19회 감정평가사 1차시험 기출문제/1교시.hwp
%% conversion: kordoc 4.24.51 -> tools/gen1.py (auto); review=proofread-sonnet; flags=none
\begin{question}{22}
물권적 청구권에 관한 설명으로 틀린 것은? (다툼이 있는 경우에는 판례에 의함)
\begin{choices}
  \item 불법점유를 이유로 한 건물명도청구를 하려면 현실적으로 불법점유하고 있는 사람을 상대로 하여야 하지만, 그렇지 않는 경우에는 간접점유자를 상대로 명도를 청구할 수 있다.
  \item 사기에 의해 물건을 인도한 자는 점유회수의 소를 제기할 수 없다.
  \item 소유권을 양도한 전소유자는 제3자인 불법점유자에 대하여 소유권에 기한 물권적 청구권에 의한 방해배제를 구할 수 있다.
  \item 점유물방해제거청구권은 방해가 종료한 날로부터 1년 내에 행사하여야 하는데, 그 기간은 출소기간으로 해석된다.
  \item 인근의 소음으로 생활이익이 침해된 건물의 소유자는 그 침해가 사회통념상 수인한도를 넘는 경우에 소유권에 기하여 소음피해의 제거청구나 예방을 위한 유지청구를 할 수 있다.
\end{choices}
\end{question}
\answer{3}
